% Options for packages loaded elsewhere
\PassOptionsToPackage{unicode}{hyperref}
\PassOptionsToPackage{hyphens}{url}
\documentclass[
]{article}
\usepackage{xcolor}
\usepackage{amsmath,amssymb}
\setcounter{secnumdepth}{-\maxdimen} % remove section numbering
\usepackage{iftex}
\ifPDFTeX
  \usepackage[T1]{fontenc}
  \usepackage[utf8]{inputenc}
  \usepackage{textcomp} % provide euro and other symbols
\else % if luatex or xetex
  \usepackage{unicode-math} % this also loads fontspec
  \defaultfontfeatures{Scale=MatchLowercase}
  \defaultfontfeatures[\rmfamily]{Ligatures=TeX,Scale=1}
\fi
\usepackage{lmodern}
\ifPDFTeX\else
  % xetex/luatex font selection
\fi
% Use upquote if available, for straight quotes in verbatim environments
\IfFileExists{upquote.sty}{\usepackage{upquote}}{}
\IfFileExists{microtype.sty}{% use microtype if available
  \usepackage[]{microtype}
  \UseMicrotypeSet[protrusion]{basicmath} % disable protrusion for tt fonts
}{}
\makeatletter
\@ifundefined{KOMAClassName}{% if non-KOMA class
  \IfFileExists{parskip.sty}{%
    \usepackage{parskip}
  }{% else
    \setlength{\parindent}{0pt}
    \setlength{\parskip}{6pt plus 2pt minus 1pt}}
}{% if KOMA class
  \KOMAoptions{parskip=half}}
\makeatother
\usepackage{color}
\usepackage{fancyvrb}
\newcommand{\VerbBar}{|}
\newcommand{\VERB}{\Verb[commandchars=\\\{\}]}
\DefineVerbatimEnvironment{Highlighting}{Verbatim}{commandchars=\\\{\}}
% Add ',fontsize=\small' for more characters per line
\newenvironment{Shaded}{}{}
\newcommand{\AlertTok}[1]{\textcolor[rgb]{1.00,0.00,0.00}{\textbf{#1}}}
\newcommand{\AnnotationTok}[1]{\textcolor[rgb]{0.38,0.63,0.69}{\textbf{\textit{#1}}}}
\newcommand{\AttributeTok}[1]{\textcolor[rgb]{0.49,0.56,0.16}{#1}}
\newcommand{\BaseNTok}[1]{\textcolor[rgb]{0.25,0.63,0.44}{#1}}
\newcommand{\BuiltInTok}[1]{\textcolor[rgb]{0.00,0.50,0.00}{#1}}
\newcommand{\CharTok}[1]{\textcolor[rgb]{0.25,0.44,0.63}{#1}}
\newcommand{\CommentTok}[1]{\textcolor[rgb]{0.38,0.63,0.69}{\textit{#1}}}
\newcommand{\CommentVarTok}[1]{\textcolor[rgb]{0.38,0.63,0.69}{\textbf{\textit{#1}}}}
\newcommand{\ConstantTok}[1]{\textcolor[rgb]{0.53,0.00,0.00}{#1}}
\newcommand{\ControlFlowTok}[1]{\textcolor[rgb]{0.00,0.44,0.13}{\textbf{#1}}}
\newcommand{\DataTypeTok}[1]{\textcolor[rgb]{0.56,0.13,0.00}{#1}}
\newcommand{\DecValTok}[1]{\textcolor[rgb]{0.25,0.63,0.44}{#1}}
\newcommand{\DocumentationTok}[1]{\textcolor[rgb]{0.73,0.13,0.13}{\textit{#1}}}
\newcommand{\ErrorTok}[1]{\textcolor[rgb]{1.00,0.00,0.00}{\textbf{#1}}}
\newcommand{\ExtensionTok}[1]{#1}
\newcommand{\FloatTok}[1]{\textcolor[rgb]{0.25,0.63,0.44}{#1}}
\newcommand{\FunctionTok}[1]{\textcolor[rgb]{0.02,0.16,0.49}{#1}}
\newcommand{\ImportTok}[1]{\textcolor[rgb]{0.00,0.50,0.00}{\textbf{#1}}}
\newcommand{\InformationTok}[1]{\textcolor[rgb]{0.38,0.63,0.69}{\textbf{\textit{#1}}}}
\newcommand{\KeywordTok}[1]{\textcolor[rgb]{0.00,0.44,0.13}{\textbf{#1}}}
\newcommand{\NormalTok}[1]{#1}
\newcommand{\OperatorTok}[1]{\textcolor[rgb]{0.40,0.40,0.40}{#1}}
\newcommand{\OtherTok}[1]{\textcolor[rgb]{0.00,0.44,0.13}{#1}}
\newcommand{\PreprocessorTok}[1]{\textcolor[rgb]{0.74,0.48,0.00}{#1}}
\newcommand{\RegionMarkerTok}[1]{#1}
\newcommand{\SpecialCharTok}[1]{\textcolor[rgb]{0.25,0.44,0.63}{#1}}
\newcommand{\SpecialStringTok}[1]{\textcolor[rgb]{0.73,0.40,0.53}{#1}}
\newcommand{\StringTok}[1]{\textcolor[rgb]{0.25,0.44,0.63}{#1}}
\newcommand{\VariableTok}[1]{\textcolor[rgb]{0.10,0.09,0.49}{#1}}
\newcommand{\VerbatimStringTok}[1]{\textcolor[rgb]{0.25,0.44,0.63}{#1}}
\newcommand{\WarningTok}[1]{\textcolor[rgb]{0.38,0.63,0.69}{\textbf{\textit{#1}}}}
\usepackage{longtable,booktabs,array}
\newcounter{none} % for unnumbered tables
\usepackage{calc} % for calculating minipage widths
% Correct order of tables after \paragraph or \subparagraph
\usepackage{etoolbox}
\makeatletter
\patchcmd\longtable{\par}{\if@noskipsec\mbox{}\fi\par}{}{}
\makeatother
% Allow footnotes in longtable head/foot
\IfFileExists{footnotehyper.sty}{\usepackage{footnotehyper}}{\usepackage{footnote}}
\makesavenoteenv{longtable}
\setlength{\emergencystretch}{3em} % prevent overfull lines
\providecommand{\tightlist}{%
  \setlength{\itemsep}{0pt}\setlength{\parskip}{0pt}}
\usepackage{bookmark}
\IfFileExists{xurl.sty}{\usepackage{xurl}}{} % add URL line breaks if available
\urlstyle{same}
% fallback for those not using the hyperref driver hyperxmp:
\makeatletter
\@ifundefined{xmpquote}{\newcommand{\xmpquote}[1]{#1}}{}
\makeatother
\hypersetup{
  hidelinks,
  pdfcreator={LaTeX via pandoc}}

\author{}
\date{}

\begin{document}

\section{Dyno Test Procedure --- CubeMars AKE80-8 with Inertia Leg \&
LocoMuJoCo Policy
Replay}\label{dyno-test-procedure-cubemars-ake80-8-with-inertia-leg-locomujoco-policy-replay}

\textbf{Phase 2 \textbar{} Status: Draft \textbar{} Date: 2026-07-07}
\textbf{Related KANs:} KAN-56 (comparative test framework) \textbf{Motor
Under Test:} CubeMars AKE80-8 KV30 (QDD actuator, 8:1 planetary)
\textbf{Board Under Test:} 96V GaN inverter (KAN-54) evaluated at both
48V (Si-equivalent baseline) and 96V.

\begin{center}\rule{0.5\linewidth}{0.5pt}\end{center}

\subsection{1. Objectives \& Thesis
Contribution}\label{objectives-thesis-contribution}

\subsubsection{1.1 Purpose}\label{purpose}

This procedure defines the comprehensive \textbf{hardware-in-the-loop
validation} of the 96V GaN drive hypothesis under realistic humanoid
gait loading. It builds upon the static \((T, \omega)\) sweep
methodology from \href{../../../docs/KAN-56_test_plan.md}{KAN-56} by
first establishing a full steady-state performance baseline, and then
extending the testing to subject the motor and inverter to
\textbf{dynamic, time-varying torque/speed profiles} extracted directly
from trained LocoMuJoCo locomotion policies. This bridges the gap
between:

\begin{itemize}
\tightlist
\item
  \textbf{Simulation} --- where we have full-body dynamics, trained
  policies, and predicted joint trajectories (from
  \texttt{MuJoCo/policies/Baseline\_Active\_Policy/})
\item
  \textbf{Hardware} --- where we measure real electrical losses, real
  thermal behaviour, and real switching waveforms on an actual motor
  driven by a real inverter
\end{itemize}

\subsubsection{1.2 Thesis Claims
Supported}\label{thesis-claims-supported}

{\def\LTcaptype{none} % do not increment counter
\begin{longtable}[]{@{}
  >{\raggedright\arraybackslash}p{(\linewidth - 2\tabcolsep) * \real{0.1842}}
  >{\raggedright\arraybackslash}p{(\linewidth - 2\tabcolsep) * \real{0.8158}}@{}}
\toprule\noalign{}
\begin{minipage}[b]{\linewidth}\raggedright
Claim
\end{minipage} & \begin{minipage}[b]{\linewidth}\raggedright
Test(s) That Provide Evidence
\end{minipage} \\
\midrule\noalign{}
\endhead
\bottomrule\noalign{}
\endlastfoot
96V GaN achieves higher cycle-averaged efficiency than a 48V
Si-equivalent baseline under gait loading & T5, T6 \\
GaN's lower switching losses enable higher \(f_{sw}\) and smaller dead
time, improving Z-width (transparent impedance) & T2, T3 \\
Higher bus voltage reduces \(I_{RMS}\) and thermal stress for the same
mechanical output & T1, T4 \\
Policy-driven load profiles produce different loss distributions than
static grids & T1 vs T5/T6 \\
Passive (activation-modulated) policies reduce electrical CoT versus
active baselines & T8 \\
GaN zero-reverse-recovery enables faster regen braking during dynamic
motions & T7 \\
\end{longtable}
}

\subsubsection{1.3 Why the AKE80-8?}\label{why-the-ake80-8}

The CubeMars AKE80-8 is a commercially available QDD actuator with an
8:1 planetary gearbox. Its specifications bracket the KBot leg
actuators:

{\def\LTcaptype{none} % do not increment counter
\begin{longtable}[]{@{}llll@{}}
\toprule\noalign{}
Parameter & AKE80-8 & KBot Hip/Knee (RS-04) & KBot Ankle (RS-02) \\
\midrule\noalign{}
\endhead
\bottomrule\noalign{}
\endlastfoot
Gear ratio & 8:1 & 9:1 & 7.75:1 \\
Peak torque & 30 Nm & 120 Nm & 17 Nm \\
\(K_t\) & 0.32 Nm/A & 2.1 Nm/A & 1.22 Nm/A \\
Rated voltage & 48 V & --- & --- \\
Pole pairs & 21 & --- & --- \\
\end{longtable}
}

The AKE80-8 can serve as a \textbf{surrogate joint actuator} for dyno
testing: its gear ratio and QDD architecture are representative of
humanoid leg joints, while its smaller torque rating keeps the dyno
power requirements manageable. All results are expressed in per-unit or
normalised form so they transfer to the full-scale RS-04/RS-02 drives.

\begin{quote}
\textbf{96V operation note:} The AKE80-8 is rated at 48V nominal. At
96V, the motor sees 2× voltage headroom, which (a) doubles the
theoretical no-load speed ceiling from 195 to \textasciitilde390 rpm,
(b) halves the phase current for the same shaft power
(\(P = V \cdot I\)), reducing \(I^2 R\) winding losses by up to 4×, and
(c) gives the FOC current loop more voltage margin for faster transient
response. The motor winding insulation must be verified for 96V
operation before testing.
\end{quote}

\begin{center}\rule{0.5\linewidth}{0.5pt}\end{center}

\subsection{2. AKE80-8 Motor
Characterization}\label{ake80-8-motor-characterization}

\subsubsection{2.1 Datasheet Parameters}\label{datasheet-parameters}

{\def\LTcaptype{none} % do not increment counter
\begin{longtable}[]{@{}lllll@{}}
\toprule\noalign{}
Parameter & Symbol & Value & Unit & Source \\
\midrule\noalign{}
\endhead
\bottomrule\noalign{}
\endlastfoot
Torque constant & \(K_t\) & 0.32 & Nm/A & Datasheet \\
Back-EMF constant & \(K_e\) & 33 & V/krpm & Datasheet \\
Phase-to-phase resistance & \(R_{pp}\) & 870 & mΩ & Datasheet \\
Phase resistance (star) & \(R_{ph}\) & 435 & mΩ & \(R_{pp}/2\) \\
Phase-to-phase inductance & \(L_{pp}\) & 990 & µH & Datasheet \\
Phase inductance (star) & \(L_{ph}\) & 495 & µH & \(L_{pp}/2\) \\
Motor constant & \(K_m\) & 0.34 & Nm/√W & Datasheet \\
Pole pairs & \(p\) & 21 & --- & Datasheet \\
Gear ratio & \(N\) & 8 & :1 & Datasheet \\
Peak output torque & \(\tau_{peak}\) & 30 & Nm & Datasheet \\
Rated output torque & \(\tau_{rated}\) & 12 & Nm & Datasheet \\
Peak current & \(I_{peak}\) & 12 & A & Datasheet \\
Rated current & \(I_{rated}\) & 4.8 & A & Datasheet \\
No-load speed & \(\omega_{NL}\) & 195 & rpm & Datasheet (at 48V) \\
Rated speed & \(\omega_{rated}\) & 150 & rpm & Datasheet \\
Electrical time constant & \(\tau_e\) & 1.13 & ms & Datasheet \\
Backlash & --- & 9 & arcmin & Datasheet \\
Weight & --- & 570 & g & Datasheet \\
\end{longtable}
}

\subsubsection{2.2 Pre-Test Measurements (To Be Filled
In)}\label{pre-test-measurements-to-be-filled-in}

Before any dynamic testing, measure and record these on the bench with
the motor disconnected from the dyno:

{\def\LTcaptype{none} % do not increment counter
\begin{longtable}[]{@{}
  >{\raggedright\arraybackslash}p{(\linewidth - 6\tabcolsep) * \real{0.2558}}
  >{\raggedright\arraybackslash}p{(\linewidth - 6\tabcolsep) * \real{0.3488}}
  >{\raggedright\arraybackslash}p{(\linewidth - 6\tabcolsep) * \real{0.2558}}
  >{\raggedright\arraybackslash}p{(\linewidth - 6\tabcolsep) * \real{0.1395}}@{}}
\toprule\noalign{}
\begin{minipage}[b]{\linewidth}\raggedright
Parameter
\end{minipage} & \begin{minipage}[b]{\linewidth}\raggedright
Measured Value
\end{minipage} & \begin{minipage}[b]{\linewidth}\raggedright
Instrument
\end{minipage} & \begin{minipage}[b]{\linewidth}\raggedright
Date
\end{minipage} \\
\midrule\noalign{}
\endhead
\bottomrule\noalign{}
\endlastfoot
\(R_{pp}\) at 25°C & \_\_\_\_\_\_\_ mΩ & LCR meter / 4-wire DMM & \\
\(L_{pp}\) at 1 kHz & \_\_\_\_\_\_\_ µH & LCR meter & \\
\(K_e\) (spin at known RPM, measure open-circuit Vpp) & \_\_\_\_\_\_\_
V/krpm & Scope + tachometer & \\
Cogging torque peak & \_\_\_\_\_\_\_ Nm & Torque sensor, slow manual
rotation & \\
Rotor inertia \(J_m\) (optional) & \_\_\_\_\_\_\_ kg·m² & Spin-down
deceleration test & \\
\end{longtable}
}

\subsubsection{2.3 Torque--Speed Envelope}\label{torquespeed-envelope}

The operating envelope shifts dramatically between 48V and 96V supply:

\begin{verbatim}
Torque (Nm)
  30 ┤ ████████                          Peak (transient < 1s)
     │ ████████
  12 ┤ ████████████████████               Rated (continuous)
     │ █████████████████████████
   0 ┤──────────────────────────────────── Speed (rpm)
     0    50   100   150   195  ~390
         ◄── 48V ──►  ◄─── 96V ───►
\end{verbatim}

At 96V, the constant-torque region extends further because the inverter
has more bus voltage to overcome the back-EMF at high speed:

\[
\omega_{max}(V_{bus}) = \frac{V_{bus} - I \cdot R_{ph}}{K_e} \approx \frac{V_{bus}}{K_e}
\]

This is the primary electrical advantage being tested: for a given gait
trajectory \((q(t), \dot{q}(t), \tau(t))\), the 96V drive operates
deeper inside its constant-torque region, where the current loop has
more voltage margin and the efficiency is highest.

\begin{center}\rule{0.5\linewidth}{0.5pt}\end{center}

\subsection{3. Inertia Leg Design --- Pendulum Arm
Fixture}\label{inertia-leg-design-pendulum-arm-fixture}

\subsubsection{3.1 Concept}\label{concept}

The test stand uses the simplest possible layout: the AKE80-8 motor is
\textbf{clamped to the edge of a sturdy table} with its output shaft
pointing horizontally. A rigid bar (the ``leg'') is bolted to the output
shaft and hangs downward under gravity, free to swing in the vertical
plane.

\begin{verbatim}
    ┌─────────────────────────┐
    │        TABLE            │
    │                         │
    │   ┌──────────┐          │
    │   │ AKE80-8  │◄── Clamped to table edge
    │   │  Motor   │
    └───┤          ├──────────┘
        └────┬─────┘
             │ Output shaft (horizontal)
             │
        ┌────┴────┐
        │ Coupler │
        └────┬────┘
             │
             │  ◄── Steel / aluminium bar ("leg")
             │
             │       Length L, adjustable masses
             │
             ●  ◄── Clamp-on mass (m_tip)
             │
             ▼ gravity
\end{verbatim}

This approach has three major advantages over a traditional
flywheel-on-horizontal-shaft dyno:

\begin{enumerate}
\def\labelenumi{\arabic{enumi}.}
\tightlist
\item
  \textbf{Free gravity loading} --- the pendulum arm experiences real
  configuration-dependent gravitational torque
  \(G(q) = -m g L_{COM} \sin(q)\), exactly like a real leg joint. No
  software feedforward or load motor is needed to emulate gravity.
\item
  \textbf{Physically intuitive} --- it literally looks like a swinging
  leg segment, making the test setup easy to explain in the thesis and
  to visitors.
\item
  \textbf{Simple construction} --- a bar, a coupler, and some clamp-on
  masses. No precision machining of flywheel discs required.
\end{enumerate}

\subsubsection{3.2 Inertia Targets}\label{inertia-targets}

The KBot's sagittal-plane leg model (from
\href{../../MuJoCo/kbot/leg_params.py}{\texttt{leg\_params.py}}) defines
the inertia the motor should ``see'' at the joint side (after the
gearbox). Three configurations:

\paragraph{Configuration A --- Knee Joint
Equivalent}\label{configuration-a-knee-joint-equivalent}

From \href{../../MuJoCo/kbot/leg_params.py}{\texttt{leg\_params.py}}
Model A (hip locked, ankle locked):

{\def\LTcaptype{none} % do not increment counter
\begin{longtable}[]{@{}
  >{\raggedright\arraybackslash}p{(\linewidth - 6\tabcolsep) * \real{0.3030}}
  >{\raggedright\arraybackslash}p{(\linewidth - 6\tabcolsep) * \real{0.2424}}
  >{\raggedright\arraybackslash}p{(\linewidth - 6\tabcolsep) * \real{0.2121}}
  >{\raggedright\arraybackslash}p{(\linewidth - 6\tabcolsep) * \real{0.2424}}@{}}
\toprule\noalign{}
\begin{minipage}[b]{\linewidth}\raggedright
Quantity
\end{minipage} & \begin{minipage}[b]{\linewidth}\raggedright
Symbol
\end{minipage} & \begin{minipage}[b]{\linewidth}\raggedright
Value
\end{minipage} & \begin{minipage}[b]{\linewidth}\raggedright
Source
\end{minipage} \\
\midrule\noalign{}
\endhead
\bottomrule\noalign{}
\endlastfoot
Shank + foot inertia about knee & \(I_{distal}\) & 0.1000 kg·m² &
\texttt{leg\_params.py} L89 \\
Reflected motor inertia at joint & \(J_{reflected}\) & 0.04 kg·m² &
\(N^2 J_m\) (RS-04) \\
\textbf{Total joint-side inertia} & \(I_{total}\) & \textbf{0.1400
kg·m²} & \texttt{leg\_params.py} L95 \\
Distal mass & \(m_{distal}\) & 2.294 kg & \(m_2 + m_3\) \\
Distal COM distance from knee & \(d_{COM}\) & 0.0966 m &
\texttt{leg\_params.py} L102 \\
Gravity torque constants & \(a_g, b_g\) & 0.069, −0.371 kg·m &
\texttt{leg\_params.py} L107-108 \\
\end{longtable}
}

\paragraph{Configuration B --- Hip Pitch
Equivalent}\label{configuration-b-hip-pitch-equivalent}

{\def\LTcaptype{none} % do not increment counter
\begin{longtable}[]{@{}
  >{\raggedright\arraybackslash}p{(\linewidth - 6\tabcolsep) * \real{0.2703}}
  >{\raggedright\arraybackslash}p{(\linewidth - 6\tabcolsep) * \real{0.2162}}
  >{\raggedright\arraybackslash}p{(\linewidth - 6\tabcolsep) * \real{0.1892}}
  >{\raggedright\arraybackslash}p{(\linewidth - 6\tabcolsep) * \real{0.3243}}@{}}
\toprule\noalign{}
\begin{minipage}[b]{\linewidth}\raggedright
Quantity
\end{minipage} & \begin{minipage}[b]{\linewidth}\raggedright
Symbol
\end{minipage} & \begin{minipage}[b]{\linewidth}\raggedright
Value
\end{minipage} & \begin{minipage}[b]{\linewidth}\raggedright
Derivation
\end{minipage} \\
\midrule\noalign{}
\endhead
\bottomrule\noalign{}
\endlastfoot
Thigh inertia at COM & \(I_1\) & 0.1272 kg·m² & \texttt{leg\_params.py}
L41 \\
Thigh mass & \(m_1\) & 5.297 kg & \texttt{leg\_params.py} L25 \\
Thigh COM distance & \(\|r_1\|\) & 0.1944 m & \texttt{leg\_params.py}
L34 \\
Thigh parallel-axis at hip & \(I_{1,hip}\) & \(I_1 + m_1 \|r_1\|^2\) =
0.327 kg·m² & Parallel axis theorem \\
Shank+foot about hip (approx.) & \(I_{2+3,hip}\) & ≈ 0.441 kg·m² &
\(I_{distal} + m_{distal} \cdot l_1^2\) \\
Reflected motor inertia & \(J_{reflected}\) & 0.04 kg·m² & \(N^2 J_m\)
(RS-04) \\
\textbf{Total joint-side inertia} & \(I_{hip,total}\) & \textbf{≈ 0.808
kg·m²} & Sum \\
Total distal mass & \(m_{leg}\) & 7.590 kg & \(m_1 + m_2 + m_3\) \\
\end{longtable}
}

\paragraph{Configuration C --- Ankle Joint
Equivalent}\label{configuration-c-ankle-joint-equivalent}

{\def\LTcaptype{none} % do not increment counter
\begin{longtable}[]{@{}
  >{\raggedright\arraybackslash}p{(\linewidth - 6\tabcolsep) * \real{0.3030}}
  >{\raggedright\arraybackslash}p{(\linewidth - 6\tabcolsep) * \real{0.2424}}
  >{\raggedright\arraybackslash}p{(\linewidth - 6\tabcolsep) * \real{0.2121}}
  >{\raggedright\arraybackslash}p{(\linewidth - 6\tabcolsep) * \real{0.2424}}@{}}
\toprule\noalign{}
\begin{minipage}[b]{\linewidth}\raggedright
Quantity
\end{minipage} & \begin{minipage}[b]{\linewidth}\raggedright
Symbol
\end{minipage} & \begin{minipage}[b]{\linewidth}\raggedright
Value
\end{minipage} & \begin{minipage}[b]{\linewidth}\raggedright
Source
\end{minipage} \\
\midrule\noalign{}
\endhead
\bottomrule\noalign{}
\endlastfoot
Foot inertia at COM & \(I_3\) & 0.00189 kg·m² & \texttt{leg\_params.py}
L43 \\
Foot mass & \(m_3\) & 0.609 kg & \texttt{leg\_params.py} L27 \\
Foot COM distance & \(\|r_3\|\) & 0.0353 m & \texttt{leg\_params.py}
L36 \\
Foot parallel-axis at ankle & \(I_{3,ankle}\) & \(I_3 + m_3 \|r_3\|^2\)
= 0.00265 kg·m² & \\
Reflected motor inertia & \(J_{reflected}\) & 0.0042 kg·m² & \(N^2 J_m\)
(RS-02) \\
\textbf{Total joint-side inertia} & \(I_{ankle,total}\) & \textbf{≈
0.0069 kg·m²} & Sum \\
\end{longtable}
}

\subsubsection{3.3 Pendulum Arm Sizing}\label{pendulum-arm-sizing}

The pendulum arm must match \textbf{two quantities simultaneously}: the
rotational inertia \(I_{target}\) and the gravity torque magnitude
\(m g L_{COM}\) of the real leg segment.

For a simple arm consisting of a light rigid bar (mass negligible) with
a concentrated tip mass \(m_{tip}\) at distance \(L\) from the pivot:

\[
I_{arm} = m_{tip} \cdot L^2
\] \[
G_{arm}(q) = -m_{tip} \cdot g \cdot L \cdot \sin(q)
\]

Matching to the KBot values requires: \[
m_{tip} = \frac{I_{target}}{L^2}, \qquad L = \frac{I_{target}}{m_{distal} \cdot d_{COM}}
\]

where \(m_{distal} \cdot d_{COM}\) is the first mass moment that sets
the gravity torque amplitude. If the bar itself has significant mass
\(m_{bar}\) distributed uniformly over length \(L_{bar}\), add
\(\frac{1}{3} m_{bar} L_{bar}^2\) to the inertia.

\paragraph{Pendulum Arm Sizing Table}\label{pendulum-arm-sizing-table}

{\def\LTcaptype{none} % do not increment counter
\begin{longtable}[]{@{}
  >{\raggedright\arraybackslash}p{(\linewidth - 10\tabcolsep) * \real{0.0870}}
  >{\raggedright\arraybackslash}p{(\linewidth - 10\tabcolsep) * \real{0.2283}}
  >{\raggedright\arraybackslash}p{(\linewidth - 10\tabcolsep) * \real{0.1739}}
  >{\raggedright\arraybackslash}p{(\linewidth - 10\tabcolsep) * \real{0.2174}}
  >{\raggedright\arraybackslash}p{(\linewidth - 10\tabcolsep) * \real{0.2174}}
  >{\raggedright\arraybackslash}p{(\linewidth - 10\tabcolsep) * \real{0.0761}}@{}}
\toprule\noalign{}
\begin{minipage}[b]{\linewidth}\raggedright
Config
\end{minipage} & \begin{minipage}[b]{\linewidth}\raggedright
\(I_{target}\) (kg·m²)
\end{minipage} & \begin{minipage}[b]{\linewidth}\raggedright
Bar Length \(L\)
\end{minipage} & \begin{minipage}[b]{\linewidth}\raggedright
Tip Mass \(m_{tip}\)
\end{minipage} & \begin{minipage}[b]{\linewidth}\raggedright
Gravity Torque Peak
\end{minipage} & \begin{minipage}[b]{\linewidth}\raggedright
Notes
\end{minipage} \\
\midrule\noalign{}
\endhead
\bottomrule\noalign{}
\endlastfoot
\textbf{A --- Knee} & 0.140 & 0.30 m & 1.56 kg & 4.58 Nm & Matches KBot
shank length \\
\textbf{A --- Knee} & 0.140 & 0.40 m & 0.88 kg & 3.44 Nm & Longer bar,
lighter mass \\
\textbf{B --- Hip} & 0.808 & 0.40 m & 5.05 kg & 19.8 Nm & Heavy --- use
thick steel bar \\
\textbf{B --- Hip} & 0.808 & 0.50 m & 3.23 kg & 15.8 Nm & More
manageable mass \\
\textbf{C --- Ankle} & 0.0069 & 0.10 m & 0.69 kg & 0.68 Nm & Short stub
+ small mass \\
\end{longtable}
}

\begin{quote}
\textbf{Recommended starting configuration:} Config A with \(L = 0.30\)
m (matches the KBot shank length \(l_2 = 0.292\) m) and
\(m_{tip} \approx 1.6\) kg. Use a 25 mm × 6 mm steel flat bar
(\(\approx 0.35\) kg for 300 mm length) with a 1.2 kg clamp-on mass at
the tip. The bar's own inertia
(\(\frac{1}{3} \times 0.35 \times 0.30^2 = 0.0105\) kg·m²) is small
relative to the tip contribution (\(1.2 \times 0.30^2 = 0.108\) kg·m²),
totalling \(\approx 0.119\) kg·m². Adjust tip mass position along the
bar to fine-tune to the exact 0.140 kg·m² target.
\end{quote}

\subsubsection{3.4 Gravity Loading --- Built
In}\label{gravity-loading-built-in}

Unlike a horizontal-shaft flywheel dyno, the pendulum arm provides
\textbf{real gravitational loading} automatically. The gravity torque at
joint angle \(q\) (measured from the vertical hang-down position) is:

\[
G_{pend}(q) = -m_{tip} \cdot g \cdot L \cdot \sin(q)
\]

For Config A (\(m_{tip} = 1.56\) kg, \(L = 0.30\) m): - Peak gravity
torque: \(1.56 \times 9.81 \times 0.30 = 4.59\) Nm - Compare to KBot
knee gravity peak:
\(g \cdot \sqrt{a_g^2 + b_g^2} = 9.81 \times 0.377 = 3.70\) Nm

The pendulum slightly over-estimates the KBot knee gravity torque (4.59
vs 3.70 Nm) because the KBot's COM is offset forward, not purely below
the joint. This is an acceptable approximation for drive-system testing
--- the motor sees realistic gravity-like loading, and the exact
magnitude difference is within 25\%.

\begin{quote}
\textbf{For the non-policy tests (T1--T4 in §5), gravity loading simply
adds realism.} The motor swings the pendulum back and forth while we
measure efficiency. For later policy-replay tests, the gravity torque
naturally shows up in the measured shaft torque.
\end{quote}

\subsubsection{3.5 Fixture Bill of
Materials}\label{fixture-bill-of-materials}

{\def\LTcaptype{none} % do not increment counter
\begin{longtable}[]{@{}
  >{\raggedright\arraybackslash}p{(\linewidth - 6\tabcolsep) * \real{0.1818}}
  >{\raggedright\arraybackslash}p{(\linewidth - 6\tabcolsep) * \real{0.4242}}
  >{\raggedright\arraybackslash}p{(\linewidth - 6\tabcolsep) * \real{0.1515}}
  >{\raggedright\arraybackslash}p{(\linewidth - 6\tabcolsep) * \real{0.2424}}@{}}
\toprule\noalign{}
\begin{minipage}[b]{\linewidth}\raggedright
Item
\end{minipage} & \begin{minipage}[b]{\linewidth}\raggedright
Specification
\end{minipage} & \begin{minipage}[b]{\linewidth}\raggedright
Qty
\end{minipage} & \begin{minipage}[b]{\linewidth}\raggedright
Purpose
\end{minipage} \\
\midrule\noalign{}
\endhead
\bottomrule\noalign{}
\endlastfoot
Table clamp / L-bracket & Steel, ≥ M8 bolts, rated ≥ 50 Nm reaction & 1
& Secure motor to table edge \\
Shaft coupler & Rigid jaw coupler, Ø8 mm bore (match AKE80-8 shaft) & 1
& Connect motor to arm \\
Steel flat bar & 25 × 6 mm, 1045 steel, 300--500 mm length & 1 &
Pendulum arm \\
Clamp-on masses & 0.5 kg and 1.0 kg, split-collar type & 2--3 &
Adjustable inertia \\
Shaft collar / set screws & Match coupler bore & 2 & Retain masses on
bar \\
End stop bolts (optional) & M6 shoulder bolts in clamp bracket & 2 &
Limit swing angle to ±90° \\
\end{longtable}
}

\subsubsection{3.6 Range-of-Motion
Considerations}\label{range-of-motion-considerations}

The pendulum can swing freely through 360° if the table overhang is
sufficient. For most gait tests, the joint angle excursion is modest:

{\def\LTcaptype{none} % do not increment counter
\begin{longtable}[]{@{}lll@{}}
\toprule\noalign{}
Gait Mode & Typical Joint Excursion & Pendulum Swing \\
\midrule\noalign{}
\endhead
\bottomrule\noalign{}
\endlastfoot
Standing / static tests & ±5° & Negligible \\
Squat & 0--60° (knee) & Moderate \\
Walking & ±30° (hip pitch) & Moderate \\
Jump & 0--90° (knee) & Large \\
\end{longtable}
}

Ensure the table overhang allows at least ±100° of unobstructed swing.
Install mechanical end stops (rubber bumpers on shoulder bolts) at ±100°
to prevent the arm from wrapping around and hitting the table or wiring
in case of a control fault.

\begin{center}\rule{0.5\linewidth}{0.5pt}\end{center}

\subsection{4. Policy-to-Hardware
Pipeline}\label{policy-to-hardware-pipeline}

\subsubsection{4.1 Overview}\label{overview}

The pipeline converts a trained LocoMuJoCo policy into real-time motor
commands:

\begin{verbatim}
┌──────────────────┐    ┌─────────────────┐    ┌──────────────────┐    ┌──────────┐
│  Trained PPO     │    │  MuJoCo Rollout  │    │  Single-Joint    │    │  MCU     │
│  Policy (.pkl)   │───►│  q(t), q̇(t),    │───►│  Trajectory      │───►│  Inverter│
│                  │    │  τ(t) @ 50 Hz    │    │  Extraction      │    │  FOC     │
└──────────────────┘    └─────────────────┘    └──────────────────┘    └──────────┘
                                                       │
                                                       ▼
                                                ┌──────────────┐
                                                │  CSV / Binary │
                                                │  trajectory   │
                                                │  file         │
                                                └──────────────┘
\end{verbatim}

\subsubsection{4.2 Step 1 --- Policy Rollout in
Simulation}\label{step-1-policy-rollout-in-simulation}

Run each trained policy (from
\texttt{MuJoCo/policies/Baseline\_Active\_Policy/}) in the full KBot
MuJoCo model and record the complete state trajectory:

\begin{Shaded}
\begin{Highlighting}[]
\CommentTok{\# Pseudocode — run in the LocoMuJoCo training environment}
\ControlFlowTok{for}\NormalTok{ policy\_name }\KeywordTok{in}\NormalTok{ [}\StringTok{"walk\_07\_12"}\NormalTok{, }\StringTok{"run\_38\_03"}\NormalTok{, }\StringTok{"jump\_75\_01"}\NormalTok{]:}
\NormalTok{    policy }\OperatorTok{=}\NormalTok{ load\_policy(}\SpecialStringTok{f"policies/Baseline\_Active\_Policy/}\SpecialCharTok{\{}\NormalTok{policy\_name}\SpecialCharTok{\}}\SpecialStringTok{/PPOJax\_saved.pkl"}\NormalTok{)}
\NormalTok{    env }\OperatorTok{=}\NormalTok{ make\_kbot\_env(model\_path}\OperatorTok{=}\StringTok{"kbot/robot\_legs.mjcf"}\NormalTok{)}
    
\NormalTok{    trajectory }\OperatorTok{=}\NormalTok{ []}
\NormalTok{    obs }\OperatorTok{=}\NormalTok{ env.reset()}
    \ControlFlowTok{for}\NormalTok{ step }\KeywordTok{in} \BuiltInTok{range}\NormalTok{(N\_steps):}
\NormalTok{        action }\OperatorTok{=}\NormalTok{ policy(obs)}
\NormalTok{        obs, reward, done, info }\OperatorTok{=}\NormalTok{ env.step(action)}
\NormalTok{        trajectory.append(\{}
            \StringTok{"t"}\NormalTok{: step }\OperatorTok{*}\NormalTok{ env.dt,}
            \StringTok{"qpos"}\NormalTok{: env.data.qpos.copy(),        }\CommentTok{\# All joint positions}
            \StringTok{"qvel"}\NormalTok{: env.data.qvel.copy(),        }\CommentTok{\# All joint velocities  }
            \StringTok{"ctrl"}\NormalTok{: env.data.ctrl.copy(),        }\CommentTok{\# All actuator commands}
            \StringTok{"qfrc\_actuator"}\NormalTok{: env.data.qfrc\_actuator.copy()  }\CommentTok{\# Actual joint torques}
\NormalTok{        \})}
\NormalTok{    save\_trajectory(trajectory, }\SpecialStringTok{f"phase2/measurements/trajectories/}\SpecialCharTok{\{}\NormalTok{policy\_name}\SpecialCharTok{\}}\SpecialStringTok{.npz"}\NormalTok{)}
\end{Highlighting}
\end{Shaded}

Alternatively, use the existing
\href{../../MuJoCo/controllers/gait_replay.py}{\texttt{GaitReplay}}
class with pre-recorded \texttt{.npz} clips to replay existing motion
capture data through the WBC controller and record the resulting joint
torques.

\subsubsection{4.3 Step 2 --- Single-Joint Trajectory
Extraction}\label{step-2-single-joint-trajectory-extraction}

Extract the trajectory for the specific joint being tested (e.g., left
knee pitch):

\begin{Shaded}
\begin{Highlighting}[]
\CommentTok{\# Extract single{-}joint trajectory for dyno replay}
\NormalTok{joint\_idx }\OperatorTok{=} \DecValTok{18}  \CommentTok{\# dof\_left\_knee\_04 nn\_id from metadata.json}

\NormalTok{q\_joint }\OperatorTok{=}\NormalTok{ trajectory[}\StringTok{"qpos"}\NormalTok{][:, joint\_qpos\_idx]     }\CommentTok{\# Joint{-}side position [rad]}
\NormalTok{qdot\_joint }\OperatorTok{=}\NormalTok{ trajectory[}\StringTok{"qvel"}\NormalTok{][:, joint\_qvel\_idx]   }\CommentTok{\# Joint{-}side velocity [rad/s]}
\NormalTok{tau\_joint }\OperatorTok{=}\NormalTok{ trajectory[}\StringTok{"qfrc\_actuator"}\NormalTok{][:, joint\_idx] }\CommentTok{\# Joint{-}side torque [Nm]}

\CommentTok{\# Convert to motor{-}side quantities for the AKE80{-}8}
\NormalTok{N }\OperatorTok{=} \DecValTok{8}  \CommentTok{\# AKE80{-}8 gear ratio}
\NormalTok{q\_motor }\OperatorTok{=}\NormalTok{ N }\OperatorTok{*}\NormalTok{ q\_joint               }\CommentTok{\# Motor position [rad]}
\NormalTok{qdot\_motor }\OperatorTok{=}\NormalTok{ N }\OperatorTok{*}\NormalTok{ qdot\_joint          }\CommentTok{\# Motor velocity [rad/s]}
\NormalTok{tau\_motor }\OperatorTok{=}\NormalTok{ tau\_joint }\OperatorTok{/}\NormalTok{ N            }\CommentTok{\# Motor torque [Nm]}
\NormalTok{I\_motor }\OperatorTok{=}\NormalTok{ tau\_motor }\OperatorTok{/}\NormalTok{ Kt\_AKE80      }\CommentTok{\# Motor current [A]  (Kt = 0.32 Nm/A)}
\end{Highlighting}
\end{Shaded}

\subsubsection{4.4 Step 3 --- Torque
Scaling}\label{step-3-torque-scaling}

The KBot RS-04 knee actuator (\(K_t = 2.1\) Nm/A, \(N = 9\),
\(\tau_{peak} = 120\) Nm) is significantly larger than the AKE80-8.
Direct replay of RS-04 torques would saturate the AKE80-8 (30 Nm peak).
We apply per-unit normalisation:

\[
\tau_{AKE80}(t) = \tau_{policy}(t) \cdot \frac{\tau_{peak,AKE80}}{\tau_{peak,RS\text{-}04}}
= \tau_{policy}(t) \cdot \frac{30}{120} = 0.25 \cdot \tau_{policy}(t)
\]

This preserves the \textbf{shape} of the torque waveform (which
determines the loss profile) while scaling the \textbf{amplitude} to fit
within the AKE80-8's capability. All efficiency results are then
reported in per-unit form:

\[
\eta(\hat{\tau}, \hat{\omega}) \quad \text{where} \quad \hat{\tau} = \frac{\tau}{\tau_{rated}}, \quad \hat{\omega} = \frac{\omega}{\omega_{rated}}
\]

Similarly, scale velocity to respect the AKE80-8 speed limit:

\[
\omega_{AKE80}(t) = \omega_{policy}(t) \cdot \frac{\omega_{NL,AKE80}}{\omega_{NL,RS\text{-}04}}
\]

\subsubsection{4.5 Step 4 --- Gravity Compensation (Not
Required)}\label{step-4-gravity-compensation-not-required}

Because the pendulum arm fixture (§3.4) provides \textbf{real
gravitational loading}, there is no need to inject a software gravity
feedforward into the torque command. The motor naturally experiences
configuration-dependent gravity torque from the swinging arm.

The policy-derived torque command is sent directly:

\[
\tau_{cmd}(t) = \tau_{AKE80,scaled}(t)
\]

\begin{quote}
\textbf{Note:} If the pendulum arm's gravity torque magnitude differs
significantly from the KBot joint's gravity torque (see §3.4
comparison), a small correction term can optionally be added:
\(\Delta G = G_{KBot}(q) - G_{pend}(q)\), scaled appropriately. For this
first test campaign, this correction is omitted --- the 25\% difference
is acceptable for drive-system characterisation.
\end{quote}

\subsubsection{4.6 Step 5 --- Real-Time Streaming to
MCU}\label{step-5-real-time-streaming-to-mcu}

The final trajectory is streamed to the inverter MCU as a time-indexed
table:

{\def\LTcaptype{none} % do not increment counter
\begin{longtable}[]{@{}
  >{\raggedright\arraybackslash}p{(\linewidth - 8\tabcolsep) * \real{0.1410}}
  >{\raggedright\arraybackslash}p{(\linewidth - 8\tabcolsep) * \real{0.2179}}
  >{\raggedright\arraybackslash}p{(\linewidth - 8\tabcolsep) * \real{0.3205}}
  >{\raggedright\arraybackslash}p{(\linewidth - 8\tabcolsep) * \real{0.2436}}
  >{\raggedright\arraybackslash}p{(\linewidth - 8\tabcolsep) * \real{0.0769}}@{}}
\toprule\noalign{}
\begin{minipage}[b]{\linewidth}\raggedright
Time (ms)
\end{minipage} & \begin{minipage}[b]{\linewidth}\raggedright
\(q_{ref}\) (rad)
\end{minipage} & \begin{minipage}[b]{\linewidth}\raggedright
\(\dot{q}_{ref}\) (rad/s)
\end{minipage} & \begin{minipage}[b]{\linewidth}\raggedright
\(\tau_{ff}\) (Nm)
\end{minipage} & \begin{minipage}[b]{\linewidth}\raggedright
Mode
\end{minipage} \\
\midrule\noalign{}
\endhead
\bottomrule\noalign{}
\endlastfoot
0 & 0.000 & 0.000 & 0.120 & POS+FF \\
20 & 0.015 & 0.750 & 0.135 & POS+FF \\
40 & 0.058 & 1.420 & 0.189 & POS+FF \\
\ldots{} & \ldots{} & \ldots{} & \ldots{} & \ldots{} \\
\end{longtable}
}

The MCU interpolates between table entries at its internal control rate
(≥20 kHz current loop). The control mode is PD position tracking with
torque feedforward:

\[
\tau_{cmd} = K_p(q_{ref} - q) + K_d(\dot{q}_{ref} - \dot{q}) + \tau_{ff}
\]

This is identical to the controller in
\href{../../docs/kbot_leg_control_math.md}{\texttt{kbot\_leg\_control\_math.md}}
§5, ensuring consistency between simulation and hardware.

\subsubsection{4.7 Which Policies to
Replay}\label{which-policies-to-replay}

{\def\LTcaptype{none} % do not increment counter
\begin{longtable}[]{@{}
  >{\raggedright\arraybackslash}p{(\linewidth - 6\tabcolsep) * \real{0.1455}}
  >{\raggedright\arraybackslash}p{(\linewidth - 6\tabcolsep) * \real{0.1455}}
  >{\raggedright\arraybackslash}p{(\linewidth - 6\tabcolsep) * \real{0.4364}}
  >{\raggedright\arraybackslash}p{(\linewidth - 6\tabcolsep) * \real{0.2727}}@{}}
\toprule\noalign{}
\begin{minipage}[b]{\linewidth}\raggedright
Policy
\end{minipage} & \begin{minipage}[b]{\linewidth}\raggedright
Source
\end{minipage} & \begin{minipage}[b]{\linewidth}\raggedright
Joint Loading Character
\end{minipage} & \begin{minipage}[b]{\linewidth}\raggedright
Primary Tests
\end{minipage} \\
\midrule\noalign{}
\endhead
\bottomrule\noalign{}
\endlastfoot
\texttt{walk} & Active \& Passive Policies & Periodic, moderate
amplitude sinusoidal loading & T5, T8 \\
\texttt{run} & Active \& Passive Policies & High frequency periodic
loading & T8 \\
\texttt{jump} & Active \& Passive Policies & Impulsive, high peak
torque, regen braking & T7, T8 \\
\texttt{squat} & Active Policy & Slow, high-torque, quasi-static & T6 \\
\texttt{step\_in\_place} & Active Policy & Periodic with impact
transients & T6 \\
\end{longtable}
}

\begin{center}\rule{0.5\linewidth}{0.5pt}\end{center}

\subsection{5. Test Matrix}\label{test-matrix}

All tests are run on the \textbf{96V GaN inverter board} driving the
AKE80-8 motor with the pendulum arm. To demonstrate the benefits of the
GaN architecture without needing a separate Si board, we define a
\textbf{``Si-equivalent baseline''} (48V bus, 20--50 kHz \(f_{sw}\), 500
ns dead time) and compare it against the \textbf{``GaN-advanced''}
configuration (96V bus, up to 100 kHz \(f_{sw}\), \textless100 ns dead
time).

Tests are ordered so that \textbf{hardware characterisation comes first}
(T1--T4), followed by \textbf{policy-driven dynamic loading} (T5--T8).

\textbf{Test order:} T1 → T2 → T3 → T4 → T5 → T6 → T7 → T8. Allow the
motor to cool to ambient between tests.

\subsubsection{T1 --- Static Efficiency Map \& Voltage
Comparison}\label{t1-static-efficiency-map-voltage-comparison}

\textbf{Purpose:} Establish the \((T, \omega)\) efficiency grid as a
reference, and prove the bus voltage advantage (48V vs 96V) under
steady-state conditions.

{\def\LTcaptype{none} % do not increment counter
\begin{longtable}[]{@{}
  >{\raggedright\arraybackslash}p{(\linewidth - 2\tabcolsep) * \real{0.6111}}
  >{\raggedright\arraybackslash}p{(\linewidth - 2\tabcolsep) * \real{0.3889}}@{}}
\toprule\noalign{}
\begin{minipage}[b]{\linewidth}\raggedright
Parameter
\end{minipage} & \begin{minipage}[b]{\linewidth}\raggedright
Value
\end{minipage} \\
\midrule\noalign{}
\endhead
\bottomrule\noalign{}
\endlastfoot
Configurations & \textbf{48V} (Si-baseline \(f_{sw}\)/DT) vs
\textbf{96V} (GaN \(f_{sw}\)/DT) \\
Load & Pendulum arm swinging at constant amplitude (set by position
command) \\
Torques & 1, 3, 6, 9, 12 Nm (output shaft) \\
Speeds & 10, 30, 60, 100, 150 rpm (output shaft) \\
Hold time & 10 s per point (after thermal settling) \\
Logging & V\_bus, I\_bus, τ\_shaft, ω\_shaft, T\_case, T\_heatsink \\
\end{longtable}
}

\textbf{Analysis:} Generate efficiency heatmap \(\eta(T, \omega)\) for
each configuration. Compare \(I_{RMS}\) between 48V and 96V (expecting
≈50\% reduction at 96V). These maps are the fundamental dataset ---
later policy tests will overlay their operating trajectories onto these
maps.

\subsubsection{T2 --- Switching Frequency
Sweep}\label{t2-switching-frequency-sweep}

\textbf{Purpose:} Quantify the efficiency vs.~switching frequency
trade-off, showing that GaN can operate at high \(f_{sw}\) without the
severe efficiency penalty seen in Si.

{\def\LTcaptype{none} % do not increment counter
\begin{longtable}[]{@{}
  >{\raggedright\arraybackslash}p{(\linewidth - 2\tabcolsep) * \real{0.6111}}
  >{\raggedright\arraybackslash}p{(\linewidth - 2\tabcolsep) * \real{0.3889}}@{}}
\toprule\noalign{}
\begin{minipage}[b]{\linewidth}\raggedright
Parameter
\end{minipage} & \begin{minipage}[b]{\linewidth}\raggedright
Value
\end{minipage} \\
\midrule\noalign{}
\endhead
\bottomrule\noalign{}
\endlastfoot
Profile & Fixed sinusoidal oscillation (e.g.~±30° at 1 Hz on pendulum
arm) \\
\(f_{sw}\) values & 20, 50, 75, 100 kHz \\
Duration & 60 s per frequency \\
Controlled variables & Same bus voltage (96V), same motion profile, same
dead time \\
\end{longtable}
}

\textbf{Analysis:} - Plot \(\eta_{cycle}\) vs.~\(f_{sw}\) (expect a
relatively flat curve for GaN) - Current ripple \(\Delta I_{pp}\)
vs.~\(f_{sw}\) - EMI spectrum comparison at each \(f_{sw}\) (if spectrum
analyser available)

\subsubsection{T3 --- Dead Time \& Z-Width
Characterisation}\label{t3-dead-time-z-width-characterisation}

\textbf{Purpose:} Directly measure the control benefits of GaN's
switching speed. Smaller dead time reduces voltage distortion, which in
turn improves the ``Z-width'' --- the range of mechanical impedances
(from free-swinging transparency to stiff position holding) the drive
can render without instability.

{\def\LTcaptype{none} % do not increment counter
\begin{longtable}[]{@{}
  >{\raggedright\arraybackslash}p{(\linewidth - 2\tabcolsep) * \real{0.6111}}
  >{\raggedright\arraybackslash}p{(\linewidth - 2\tabcolsep) * \real{0.3889}}@{}}
\toprule\noalign{}
\begin{minipage}[b]{\linewidth}\raggedright
Parameter
\end{minipage} & \begin{minipage}[b]{\linewidth}\raggedright
Value
\end{minipage} \\
\midrule\noalign{}
\endhead
\bottomrule\noalign{}
\endlastfoot
Profile & Zero-torque command (transparency) \& high-stiffness position
hold \\
Dead times & 500 ns (Si-equivalent), 200 ns, 100 ns, 50 ns
(GaN-optimised) \\
Controlled variables & \(f_{sw}\) = 50 kHz, 96V bus \\
\end{longtable}
}

\textbf{Analysis:} - \textbf{Z-width:} Measure uncommanded drag torque
(friction) during manual backdriving at zero torque. - Voltage
distortion and phase current THD at different dead times. - Show that 50
ns dead time allows significantly higher virtual stiffness before limit
cycles (chatter) occur.

\subsubsection{T4 --- Thermal Soak}\label{t4-thermal-soak}

\textbf{Purpose:} Characterise steady-state thermal performance at
worst-case operating point (highest losses from T1 map). Determines the
continuous rating of the drive under realistic loading.

{\def\LTcaptype{none} % do not increment counter
\begin{longtable}[]{@{}
  >{\raggedright\arraybackslash}p{(\linewidth - 2\tabcolsep) * \real{0.6111}}
  >{\raggedright\arraybackslash}p{(\linewidth - 2\tabcolsep) * \real{0.3889}}@{}}
\toprule\noalign{}
\begin{minipage}[b]{\linewidth}\raggedright
Parameter
\end{minipage} & \begin{minipage}[b]{\linewidth}\raggedright
Value
\end{minipage} \\
\midrule\noalign{}
\endhead
\bottomrule\noalign{}
\endlastfoot
Configurations & \textbf{48V} (Si-baseline) vs \textbf{96V}
(GaN-advanced) \\
Profile & Continuous oscillation at the worst-case (T,ω) from T1 \\
Duration & 20--30 min or until \(dT/dt < 0.5\) °C/min \\
Logging & T\_case, T\_heatsink, T\_winding, T\_ambient @ 1 Hz;
electrical @ 10 kHz \\
\end{longtable}
}

\textbf{Analysis:} - Thermal time constant \(\tau_{th}\) from
exponential fit - Maximum continuous torque at thermal limit for each
configuration - Show that 96V operation allows significantly longer
sustained peak torque due to lower \(I^2 R\) heating.

\subsubsection{Phase A Conclusion --- Parameter
Optimization}\label{phase-a-conclusion-parameter-optimization}

Before proceeding to Phase B, the data from T1--T4 must be used to
\textbf{select the optimal FOC parameters} for the GaN drive. These
locked-in parameters will define the ``GaN-advanced'' configuration used
for all policy replay tests:

\begin{enumerate}
\def\labelenumi{\arabic{enumi}.}
\tightlist
\item
  \textbf{Optimal \(f_{sw}\):} Chosen from T2 to balance switching
  losses against current ripple.
\item
  \textbf{Optimal Dead Time:} Chosen from T3 to maximize Z-width
  (transparency) without causing shoot-through.
\item
  \textbf{FOC Tuning:} Current loop PI gains (\(K_p, K_i\)) optimized
  for the selected \(f_{sw}\) and 96V bus to ensure maximum tracking
  bandwidth.
\end{enumerate}

\begin{center}\rule{0.5\linewidth}{0.5pt}\end{center}

\subsubsection{Phase B --- Policy-Driven Dynamic
Loading}\label{phase-b-policy-driven-dynamic-loading}

\begin{quote}
\textbf{Prerequisites:} T1--T4 complete. Measurement chain validated.
Optimal GaN parameters (\(f_{sw}\), dead time, PI gains) locked in based
on Phase A results. Policy trajectory files exported (see §4). Pendulum
arm inertia verified.
\end{quote}

\subsubsection{T5 --- Walking Policy
Replay}\label{t5-walking-policy-replay}

\textbf{Purpose:} Measure cycle-averaged electrical efficiency under
realistic periodic walking loading.

{\def\LTcaptype{none} % do not increment counter
\begin{longtable}[]{@{}ll@{}}
\toprule\noalign{}
Parameter & Value \\
\midrule\noalign{}
\endhead
\bottomrule\noalign{}
\endlastfoot
Policy & \texttt{walk} policy trajectory (active) \\
Duration & 60 s continuous (≈60 gait cycles at 1 Hz stepping) \\
Pendulum arm & Config A (knee) or Config B (hip) \\
Logging rate & ≥ 10 kHz synchronised \\
\end{longtable}
}

\textbf{Analysis:} - Cycle-averaged efficiency:
\(\eta_{cycle} = \frac{\int_0^{T_{cycle}} P_{out} \, dt}{\int_0^{T_{cycle}} P_{in} \, dt}\)
- \(I_{RMS}\) per gait cycle - Thermal rise over 60 s - Loss
segregation: conduction (\(I^2 R\)) vs switching
(\(f_{sw} \cdot E_{sw}\)) - Overlay the trajectory's \((T, \omega)\)
trace on the T1 efficiency map

\subsubsection{T6 --- Squat and Step-in-Place
Replay}\label{t6-squat-and-step-in-place-replay}

\textbf{Purpose:} Characterise performance under two additional gait
primitives: - \textbf{Squat} --- high-torque, low-speed quasi-static
loading (worst case for conduction losses) - \textbf{Step-in-place} ---
periodic loading with impact-like transients

{\def\LTcaptype{none} % do not increment counter
\begin{longtable}[]{@{}
  >{\raggedright\arraybackslash}p{(\linewidth - 2\tabcolsep) * \real{0.6111}}
  >{\raggedright\arraybackslash}p{(\linewidth - 2\tabcolsep) * \real{0.3889}}@{}}
\toprule\noalign{}
\begin{minipage}[b]{\linewidth}\raggedright
Parameter
\end{minipage} & \begin{minipage}[b]{\linewidth}\raggedright
Value
\end{minipage} \\
\midrule\noalign{}
\endhead
\bottomrule\noalign{}
\endlastfoot
Policy & \texttt{squat} and \texttt{step\_in\_place} policies \\
Profile & 0.3 Hz (squat) and 1 Hz (step) \\
Duration & 30 s squat + 60 s stepping \\
Focus metrics & Peak \(\tau\), peak \(I\), bus ripple (squat); CoT,
transient quality (step) \\
\end{longtable}
}

\textbf{Analysis:} - Squat: GaN's lower \(R_{DS(on)}\) advantage under
high-current, low-speed stress - Step: Electrical Cost of Transport
\(CoT_e\), current tracking at swing/stance transitions

\subsubsection{T7 --- Jump Policy Replay}\label{t7-jump-policy-replay}

\textbf{Purpose:} Characterise peak power delivery and regenerative
braking under highly dynamic impulsive loading.

{\def\LTcaptype{none} % do not increment counter
\begin{longtable}[]{@{}
  >{\raggedright\arraybackslash}p{(\linewidth - 2\tabcolsep) * \real{0.6111}}
  >{\raggedright\arraybackslash}p{(\linewidth - 2\tabcolsep) * \real{0.3889}}@{}}
\toprule\noalign{}
\begin{minipage}[b]{\linewidth}\raggedright
Parameter
\end{minipage} & \begin{minipage}[b]{\linewidth}\raggedright
Value
\end{minipage} \\
\midrule\noalign{}
\endhead
\bottomrule\noalign{}
\endlastfoot
Profile & \texttt{jump} policy trajectory (active) \\
Duration & 2 s per jump × 5 repetitions, 30 s rest between \\
Focus metric & Peak current, peak power, regen energy, bus voltage
transient \\
\end{longtable}
}

\textbf{Analysis:} - Peak instantaneous power
\(P_{peak} = \tau_{peak} \cdot \omega_{peak}\) - Regenerative energy
captured: \(E_{regen} = \int_{braking} P_{bus} \, dt\) (when
\(I_{bus} < 0\)) - GaN advantage: zero-reverse-recovery diode loss
during regen commutation

\subsubsection{T8 --- Active vs.~Passive Policy
Comparison}\label{t8-active-vs.-passive-policy-comparison}

\textbf{Purpose:} Directly measure the electrical energy savings of the
activation-modulated passive policy (from the
\href{../../MuJoCo/docs/passive_policy_research.md}{Hybrid
Constrained-Activation framework}) versus the standard active baseline.

{\def\LTcaptype{none} % do not increment counter
\begin{longtable}[]{@{}
  >{\raggedright\arraybackslash}p{(\linewidth - 2\tabcolsep) * \real{0.6111}}
  >{\raggedright\arraybackslash}p{(\linewidth - 2\tabcolsep) * \real{0.3889}}@{}}
\toprule\noalign{}
\begin{minipage}[b]{\linewidth}\raggedright
Parameter
\end{minipage} & \begin{minipage}[b]{\linewidth}\raggedright
Value
\end{minipage} \\
\midrule\noalign{}
\endhead
\bottomrule\noalign{}
\endlastfoot
Profiles & Active vs.~Passive implementations of \texttt{walk},
\texttt{run}, and \texttt{jump} \\
Duration & 60 s each \\
Controlled variables & Same pendulum arm, same inverter, same motor,
same speed \\
\end{longtable}
}

\textbf{Analysis:} -
\(\Delta I_{RMS} = I_{RMS,active} - I_{RMS,passive}\) -
\(\Delta P_{loss} = P_{loss,active} - P_{loss,passive}\) -
\(\Delta \eta = \eta_{passive} - \eta_{active}\) - Per-joint activation
parameter \(\alpha(t)\) correlation with measured current reduction

\begin{quote}
\textbf{Note:} This test requires the passive policy to be fully
trained. If not ready by the time the dyno is operational, run T8 later
as a follow-up. Tests T1--T7 are all executable with existing hardware
and active baseline policies.
\end{quote}

\begin{center}\rule{0.5\linewidth}{0.5pt}\end{center}

\subsection{6. Instrumentation \& Measurement
Setup}\label{instrumentation-measurement-setup}

\subsubsection{6.1 Equipment List}\label{equipment-list}

{\def\LTcaptype{none} % do not increment counter
\begin{longtable}[]{@{}
  >{\raggedright\arraybackslash}p{(\linewidth - 4\tabcolsep) * \real{0.3000}}
  >{\raggedright\arraybackslash}p{(\linewidth - 4\tabcolsep) * \real{0.2250}}
  >{\raggedright\arraybackslash}p{(\linewidth - 4\tabcolsep) * \real{0.4750}}@{}}
\toprule\noalign{}
\begin{minipage}[b]{\linewidth}\raggedright
Instrument
\end{minipage} & \begin{minipage}[b]{\linewidth}\raggedright
Purpose
\end{minipage} & \begin{minipage}[b]{\linewidth}\raggedright
Min. Specification
\end{minipage} \\
\midrule\noalign{}
\endhead
\bottomrule\noalign{}
\endlastfoot
DC power supply (programmable) & Bus voltage source & 0--100 V, ≥ 15 A,
current-limited \\
Power analyser or precision shunt + DMM & DC bus power:
\(P_{in} = V_{bus} \cdot I_{bus}\) & ±0.1\% accuracy, ≥ 100 kHz BW \\
Inline torque sensor & Shaft torque \(\tau_{shaft}\) & ≥ 50 Nm range, ≤
0.1\% FS accuracy \\
Rotary encoder & Motor position \(\theta\), velocity \(\omega\) & ≥ 4096
CPR, index pulse \\
Current probes (×3) & Phase currents \(i_a, i_b, i_c\) & Hall-effect or
Rogowski, ≥ 500 kHz BW \\
Differential voltage probes (×2) & \(V_{bus}\) and \(V_{phase}\) & ≥ 100
V, ≥ 100 MHz BW \\
Thermocouples (×4) & \(T_{GaN/FET}\), \(T_{heatsink}\), \(T_{winding}\)
(if accessible), \(T_{ambient}\) & Type K, ±1°C \\
Oscilloscope & Switching waveform capture (\(V_{DS}\), \(V_{GS}\),
\(i_{phase}\)) & ≥ 200 MHz, ≥ 1 GS/s \\
DAQ system & Synchronised multi-channel logging & ≥ 10 kHz sample rate,
≥ 8 channels \\
LISN (optional) & Conducted EMI measurement & Per CISPR standards \\
Spectrum analyser (optional) & EMI frequency analysis & 150 kHz -- 30
MHz \\
Thermal camera (optional) & Hotspot mapping & FLIR or equivalent \\
\end{longtable}
}

\subsubsection{6.2 Measurement
Architecture}\label{measurement-architecture}

\begin{verbatim}
                    ┌─────────────┐
    DC Supply ──────┤ Power       │
    (48V or 96V)    │ Analyser    ├──── V_bus, I_bus, P_in  ──► DAQ Ch 1-3
                    └──────┬──────┘
                           │ DC Bus
                    ┌──────▼──────┐
                    │  Inverter   │
                    │  Board      │──── T_case (TC) ────────► DAQ Ch 7
                    │  (GaN/Si)   │
                    └──────┬──────┘
                      3-ph │ AC    ──── i_a, i_b, i_c ─────► DAQ Ch 4-6
                    ┌──────▼──────┐         (current probes)
                    │  AKE80-8    │
                    │  Motor      │──── T_winding (TC) ─────► DAQ Ch 8
                    └──────┬──────┘
                           │ Shaft
                    ┌──────▼──────┐
                    │  Torque     │──── τ_shaft ────────────► DAQ Ch 9
                    │  Sensor     │──── ω_shaft ────────────► DAQ Ch 10
                    └──────┬──────┘
                           │
                    ┌──────▼──────┐
                    │  Pendulum   │
                    │  Arm + Mass │
                    └──────│──────┘
                           │
                           ● tip mass
                           │
                           ▼ gravity
\end{verbatim}

\subsubsection{6.3 Wiring \& Grounding
Rules}\label{wiring-grounding-rules}

\begin{enumerate}
\def\labelenumi{\arabic{enumi}.}
\tightlist
\item
  \textbf{Star ground} --- single ground reference point for all
  measurement equipment
\item
  \textbf{Identical DC bus wiring} for both inverter boards (same cable
  length, gauge, connectors)
\item
  \textbf{Shield all analogue signals} --- twisted-pair or coax from
  sensors to DAQ
\item
  \textbf{Keep current probe loops small} --- route phase wires through
  probes close to inverter output
\item
  \textbf{Oscilloscope probe ground} --- use tip-and-barrel differential
  probing for \(V_{DS}\); do NOT use scope ground clip at high-side
\end{enumerate}

\subsubsection{6.4 Safety Requirements}\label{safety-requirements}

\begin{enumerate}
\def\labelenumi{\arabic{enumi}.}
\tightlist
\item
  \textbf{96V hazard} --- while below typical ``dangerous voltage''
  thresholds, 96V can deliver lethal current through low-impedance
  paths. Use insulated connectors, shrouded terminals, and do not work
  on energised circuits.
\item
  \textbf{Emergency stop (E-stop)} --- hardware E-stop that disconnects
  DC bus. Reachable from the operator position. Tested before every
  session.
\item
  \textbf{Over-current protection} --- DC supply current limit set to 15
  A. Inverter-side over-current trip at 15 A phase (both SW and HW).
\item
  \textbf{Swing zone clearance} --- the pendulum arm sweeps a vertical
  arc. Mark a clear zone on the floor (radius = arm length + 15 cm
  margin). No hands, cables, or equipment in the swing zone during
  operation. Install mechanical end stops at ±100°.
\item
  \textbf{Thermal shutdown} --- software thermal limit at
  \(T_{case} = 100°C\). Manual abort if any component exceeds 80°C
  unexpectedly.
\item
  \textbf{Motor insulation check} --- before applying 96V, verify motor
  winding insulation with a megohmmeter (hipot) at ≥ 250V DC, confirm ≥
  100 MΩ.
\end{enumerate}

\begin{center}\rule{0.5\linewidth}{0.5pt}\end{center}

\subsection{7. Test Procedures
(Step-by-Step)}\label{test-procedures-step-by-step}

\subsubsection{7.1 Common Pre-Test
Checklist}\label{common-pre-test-checklist}

Run this before \textbf{every} test session (adapted from KAN-56 §4.1):

\begin{itemize}
\tightlist
\item[$\square$]
  Visually inspect motor, clamp, coupler, pendulum arm, torque sensor,
  and all wiring
\item[$\square$]
  Verify pendulum arm masses are secure (clamp screws tight, no play)
\item[$\square$]
  Confirm swing zone is clear of cables, tools, and personnel
\item[$\square$]
  Verify gate drive supplies (+5V for GaN, +12V for Si gate drivers)
  with no HV applied
\item[$\square$]
  Verify encoder readings: rotate shaft by hand, confirm position and
  velocity are correct
\item[$\square$]
  Verify thermocouple readings: all reading ambient ±2°C
\item[$\square$]
  Set DC supply current limit (15 A)
\item[$\square$]
  Power up at reduced bus voltage (20V), run no-load spin at low speed,
  verify:

  \begin{itemize}
  \tightlist
  \item
    Phase currents are balanced
  \item
    No audible noise, vibration, or hot spots
  \item
    Switching waveforms (\(V_{DS}\), \(V_{GS}\)) are clean on scope
  \end{itemize}
\item[$\square$]
  Ramp to full bus voltage (48V or 96V), repeat checks
\item[$\square$]
  Arm and test E-stop
\item[$\square$]
  Start DAQ logging, verify all channels are acquiring
\item[$\square$]
  Record ambient temperature and humidity in log
\item[$\square$]
  Document any deviations in \texttt{RESEARCH\_LOG.md}
\end{itemize}

\subsubsection{7.2 Per-Test Procedure
Template}\label{per-test-procedure-template}

For each test (T1--T8):

\begin{enumerate}
\def\labelenumi{\arabic{enumi}.}
\tightlist
\item
  \textbf{Load the trajectory file} for this test onto the MCU (or
  select the correct profile)
\item
  \textbf{Configure test parameters} (bus voltage, switching frequency,
  pendulum arm config)
\item
  \textbf{Start DAQ recording} --- filename format:
  \texttt{YYYY-MM-DD\_T\{N\}\_\{board\}\_\{fsw\}kHz\_\{Vbus\}V.csv}
\item
  \textbf{Start the trajectory replay} on the inverter
\item
  \textbf{Monitor} temperatures, currents, and bus voltage in real time
  during the test
\item
  \textbf{Wait} for the prescribed duration
\item
  \textbf{Stop trajectory replay} and allow motor to coast to a stop
\item
  \textbf{Stop DAQ recording}
\item
  \textbf{Allow motor to cool} to within 5°C of ambient before next test
\item
  \textbf{Transfer raw data} to
  \texttt{data\_raw/YYYY-MM-DD\_phase2\_dyno\_T\{N\}/}
\item
  \textbf{Log observations} in \texttt{RESEARCH\_LOG.md}
\end{enumerate}

\begin{center}\rule{0.5\linewidth}{0.5pt}\end{center}

\subsection{8. Data Processing \&
Analysis}\label{data-processing-analysis}

\subsubsection{8.1 Primary Metrics}\label{primary-metrics}

All analysis scripts should be placed in
\texttt{phase2/measurements/analysis/} and produce publication-ready
figures.

\paragraph{Instantaneous Efficiency}\label{instantaneous-efficiency}

\[
\eta(t) = \frac{P_{out}(t)}{P_{in}(t)} = \frac{\tau_{shaft}(t) \cdot \omega_{shaft}(t)}{V_{bus}(t) \cdot I_{bus}(t)}
\]

\paragraph{Cycle-Averaged Efficiency}\label{cycle-averaged-efficiency}

\[
\eta_{cycle} = \frac{\int_0^{T_{cycle}} \tau \cdot \omega \, dt}{\int_0^{T_{cycle}} V_{bus} \cdot I_{bus} \, dt}
= \frac{E_{mech,cycle}}{E_{elec,cycle}}
\]

\paragraph{RMS Phase Current}\label{rms-phase-current}

\[
I_{RMS} = \sqrt{\frac{1}{T} \int_0^T i_a^2(t) \, dt}
\]

\paragraph{Electrical Cost of
Transport}\label{electrical-cost-of-transport}

\[
CoT_e = \frac{E_{elec,cycle}}{m_{robot} \cdot g \cdot d_{stride}}
\]

where \(m_{robot} = 21.07\) kg (legs-only mass from
\texttt{leg\_params.py}: \(2 \times (m_1 + m_2 + m_3)
+ \text{torso}\)), \(g = 9.81\) m/s², and \(d_{stride}\) is the
displacement per gait cycle from the policy. Since we are testing a
single joint, report this as a \textbf{per-joint contribution} to the
total CoT.

\paragraph{Loss Segregation}\label{loss-segregation}

\[
P_{total} = P_{in} - P_{out}
\] \[
P_{Cu} = 3 \cdot I_{RMS}^2 \cdot R_{ph}(T_{winding}) \quad \text{(copper losses)}
\] \[
P_{Fe} \approx P_{total} - P_{Cu} - P_{sw} \quad \text{(iron/core losses, by subtraction)}
\] \[
P_{sw} \approx f_{sw} \cdot (E_{on} + E_{off}) \quad \text{(switching losses, from waveform analysis)}
\]

Note: \(R_{ph}\) must be temperature-corrected using the measured
winding temperature: \[
R_{ph}(T) = R_{ph,25°C} \cdot \left(1 + \alpha_{Cu}(T - 25°C)\right), \quad \alpha_{Cu} = 0.00393 \, °C^{-1}
\]

\subsubsection{8.2 Thermal Analysis}\label{thermal-analysis}

\begin{itemize}
\tightlist
\item
  \textbf{Thermal time constant} \(\tau_{th}\): fit an exponential
  \(T(t) = T_\infty - (T_\infty - T_0) e^{-t/\tau_{th}}\) to the soak
  test temperature data
\item
  \textbf{Thermal resistance}
  \(R_{th,j-a} = (T_{case} - T_{ambient}) / P_{loss}\) at steady state
\end{itemize}

\subsubsection{8.3 Statistical
Requirements}\label{statistical-requirements}

\begin{itemize}
\tightlist
\item
  Each test must be repeated \textbf{≥ 3 times} (minimum)
\item
  Report \textbf{mean ± standard deviation} for all primary metrics
\item
  Use the same ambient conditions (within ±3°C) for all repeats
\item
  If any repeat deviates by \textgreater{} 2σ, investigate and either
  explain or discard (with justification)
\end{itemize}

\subsubsection{8.4 Key Comparison Plots (for
Thesis)}\label{key-comparison-plots-for-thesis}

{\def\LTcaptype{none} % do not increment counter
\begin{longtable}[]{@{}
  >{\raggedright\arraybackslash}p{(\linewidth - 6\tabcolsep) * \real{0.2500}}
  >{\raggedright\arraybackslash}p{(\linewidth - 6\tabcolsep) * \real{0.2500}}
  >{\raggedright\arraybackslash}p{(\linewidth - 6\tabcolsep) * \real{0.2500}}
  >{\raggedright\arraybackslash}p{(\linewidth - 6\tabcolsep) * \real{0.2500}}@{}}
\toprule\noalign{}
\begin{minipage}[b]{\linewidth}\raggedright
Figure
\end{minipage} & \begin{minipage}[b]{\linewidth}\raggedright
X-axis
\end{minipage} & \begin{minipage}[b]{\linewidth}\raggedright
Y-axis
\end{minipage} & \begin{minipage}[b]{\linewidth}\raggedright
Curves
\end{minipage} \\
\midrule\noalign{}
\endhead
\bottomrule\noalign{}
\endlastfoot
Efficiency map & Speed (rpm) & Torque (Nm) & Colour = η; one map per
inverter \\
Walking efficiency time series & Time (s) & η(t), I(t), τ(t) & GaN vs Si
overlay \\
Cycle efficiency bar chart & Test (T1--T8) & η\_cycle (\%) & GaN vs Si
side-by-side \\
I\_RMS comparison & Gait mode & I\_RMS (A) & 48V vs 96V \\
Thermal rise & Time (s) & ΔT (°C) & GaN vs Si during T2 \\
Switching freq sweep & f\_sw (kHz) & η\_cycle (\%) & GaN and Si \\
Loss breakdown & Component & P\_loss (W) & Stacked bar: P\_Cu, P\_sw,
P\_Fe \\
Voltage comparison & Metric & Value & 48V vs 96V on same plot \\
\end{longtable}
}

\begin{center}\rule{0.5\linewidth}{0.5pt}\end{center}

\subsection{9. Expected Results \& Thesis
Integration}\label{expected-results-thesis-integration}

\subsubsection{9.1 Predicted Outcomes}\label{predicted-outcomes}

Based on the inverter specifications and motor parameters:

{\def\LTcaptype{none} % do not increment counter
\begin{longtable}[]{@{}
  >{\raggedright\arraybackslash}p{(\linewidth - 4\tabcolsep) * \real{0.1951}}
  >{\raggedright\arraybackslash}p{(\linewidth - 4\tabcolsep) * \real{0.5366}}
  >{\raggedright\arraybackslash}p{(\linewidth - 4\tabcolsep) * \real{0.2683}}@{}}
\toprule\noalign{}
\begin{minipage}[b]{\linewidth}\raggedright
Metric
\end{minipage} & \begin{minipage}[b]{\linewidth}\raggedright
Expected GaN Advantage
\end{minipage} & \begin{minipage}[b]{\linewidth}\raggedright
Reasoning
\end{minipage} \\
\midrule\noalign{}
\endhead
\bottomrule\noalign{}
\endlastfoot
Cycle η (walking) & +2--4\% absolute & Lower \(R_{DS(on)}\) + zero
reverse recovery \\
\(I_{RMS}\) (96V vs 48V) & −45--50\% & \(I \propto P/V\); halving at
doubled voltage \\
Thermal rise & −30--40\% & \(P_{Cu} \propto I^2\); 4× reduction at
96V \\
Max usable \(f_{sw}\) & 100 kHz (GaN) vs 50 kHz (Si) & GaN zero-recovery
allows aggressive timing \\
η sensitivity to \(f_{sw}\) & Flat to 100 kHz (GaN); −1\%/25 kHz (Si) &
GaN switching losses scale weakly with frequency \\
Regen braking efficiency & +5--10\% (GaN) & No diode recovery losses
during braking commutation \\
\end{longtable}
}

\subsubsection{9.2 Thesis Chapter Mapping}\label{thesis-chapter-mapping}

{\def\LTcaptype{none} % do not increment counter
\begin{longtable}[]{@{}
  >{\raggedright\arraybackslash}p{(\linewidth - 4\tabcolsep) * \real{0.1667}}
  >{\raggedright\arraybackslash}p{(\linewidth - 4\tabcolsep) * \real{0.4444}}
  >{\raggedright\arraybackslash}p{(\linewidth - 4\tabcolsep) * \real{0.3889}}@{}}
\toprule\noalign{}
\begin{minipage}[b]{\linewidth}\raggedright
Test
\end{minipage} & \begin{minipage}[b]{\linewidth}\raggedright
Chapter Section
\end{minipage} & \begin{minipage}[b]{\linewidth}\raggedright
Figure/Table
\end{minipage} \\
\midrule\noalign{}
\endhead
\bottomrule\noalign{}
\endlastfoot
T1 & Ch. 4: Motor \& Inverter Characterization & Fig. 4.x: Efficiency
map overlay \\
T2 & Ch. 4: Switching Frequency Study & Fig. 4.x: η vs f\_sw \\
T3 & Ch. 4: Voltage Architecture & Fig. 4.x: 48V vs 96V comparison \\
T4 & Ch. 4: Thermal Performance & Fig. 4.x: Thermal soak curves \\
T5 & Ch. 5: Dynamic Performance & Fig. 5.x: Walking cycle efficiency \\
T6 & Ch. 5: Dynamic Performance & Fig. 5.x: Squat \& stepping CoT \\
T7 & Ch. 5: Dynamic Performance & Fig. 5.x: Jump peak power \& regen \\
T8 & Ch. 6: Policy-Hardware Co-design & Fig. 6.x: Active vs passive
energy \\
\end{longtable}
}

\subsubsection{9.3 If Results Disagree with
Predictions}\label{if-results-disagree-with-predictions}

If the measured GaN advantage is smaller than predicted: 1. Check
thermal management --- if GaN runs hotter than expected, \(R_{DS(on)}\)
rises and advantage erodes 2. Check layout parasitics --- high loop
inductance can cause ringing and increase switching losses 3. Check dead
time --- too-conservative dead time on GaN wastes the zero-recovery
advantage 4. Revisit the loss model: are iron losses (not captured in
the conduction/switching model) dominating?

Document all discrepancies in \texttt{RESEARCH\_LOG.md} with hypotheses
and follow-up actions.

\begin{center}\rule{0.5\linewidth}{0.5pt}\end{center}

\subsection{10. File \& Data Organisation}\label{file-data-organisation}

All files produced by this procedure should follow the project
convention:

\begin{verbatim}
96v_gan_humanoid_drive/
├── phase2/
│   └── measurements/
│       ├── dyno_test_procedure_AKE80-8.md   ← This document
│       ├── trajectories/                     ← Extracted single-joint policy CSVs
│       │   ├── walk_07_12_knee.csv
│       │   ├── squat_knee.csv
│       │   ├── step_in_place_knee.csv
│       │   └── jump_75_01_knee.csv
│       └── analysis/                         ← MATLAB/Python analysis scripts
│           ├── plot_efficiency_map.py
│           ├── plot_cycle_efficiency.py
│           ├── loss_segregation.py
│           └── thermal_analysis.py
├── data_raw/
│   ├── 2026-XX-XX_phase2_dyno_T1/
│   ├── 2026-XX-XX_phase2_dyno_T2/
│   └── ...
└── notebooks/
    └── RESEARCH_LOG.md                       ← Log every test session here
\end{verbatim}

\begin{center}\rule{0.5\linewidth}{0.5pt}\end{center}

\subsection{Appendix A --- AKE80-8 Winding Insulation Verification for
96V}\label{appendix-a-ake80-8-winding-insulation-verification-for-96v}

The AKE80-8 is rated for 48V. Operating at 96V requires:

\begin{enumerate}
\def\labelenumi{\arabic{enumi}.}
\tightlist
\item
  \textbf{Phase-to-ground insulation test}: Apply 250V DC between any
  phase terminal and the motor housing using a megohmmeter. Confirm ≥
  100 MΩ insulation resistance.
\item
  \textbf{Phase-to-phase insulation test}: Apply 250V DC between any two
  phase terminals. Confirm ≥ 100 MΩ.
\item
  \textbf{Dielectric withstand (hipot)}: Apply 500V AC (or 700V DC) for
  60 seconds between phases and ground. No breakdown.
\item
  \textbf{Thermal derating}: At 96V, if the motor is driven at higher
  speed or power than the 48V rating, monitor winding temperature
  carefully. The insulation class (typically Class B or F for CubeMars
  motors) determines the maximum allowable temperature.
\end{enumerate}

\begin{quote}
\textbf{If the motor fails any insulation test, do NOT proceed with 96V
testing.} Use 48V only, or source a motor with verified higher voltage
insulation.
\end{quote}

\begin{center}\rule{0.5\linewidth}{0.5pt}\end{center}

\subsection{Appendix B --- Pendulum Arm Fixture Drawing
Reference}\label{appendix-b-pendulum-arm-fixture-drawing-reference}

A detailed technical drawing should be created in CAD (Fusion 360 or
SolidWorks) for the pendulum arm fixture assembly. Key dimensions to
specify:

\begin{itemize}
\tightlist
\item
  Table clamp geometry and bolt pattern
\item
  Shaft coupler bore diameter (match AKE80-8 output shaft, typically Ø8
  mm)
\item
  Steel flat bar cross-section, length, and material (1045 steel)
\item
  Clamp-on mass dimensions and attachment method (set screw / split
  collar)
\item
  Mass positions along bar for each inertia configuration (A, B, C)
\item
  Mechanical end stop locations (±100° from vertical)
\item
  Overall swing envelope (for safety zone marking)
\end{itemize}

Store drawing files in \texttt{phase2/mechanical/pendulum\_arm/}.

\begin{center}\rule{0.5\linewidth}{0.5pt}\end{center}

\emph{This is a living document. Update as hardware is built, tests are
executed, and results come in. Link updates to the corresponding Jira
ticket and log changes in \texttt{RESEARCH\_LOG.md}.}

\end{document}
